\scanpage{36}
% archive foliation: 35

torseurs sur $S$ de groupe $\mu_n$, et son $\pi_1$ \deleted{\unsure{est}} s'insère dans
\diagram{\begin{tikzcd}
0 \arrow[r] & \unsure{H^1(S,\mu_n)} \arrow[r] &
\pi_0(\uline{\operatorname{Brdt}}(C,\tau)) \arrow[r] &
\deleted{\unsure{\mathrm{Br}(C)}}\inserted{\operatorname{Ram}(C/A)} \arrow[r] & 0 .
\end{tikzcd}}
\deleted{(Si $\operatorname{Pic}(A)=0$, \unsure{cette catégorie}}
\[
\text{\uline{où}}\quad
\operatorname{Ram}(C/A)
\underset{\deleted{\unsure{=}}}{\overset{\mathrm{déf}}{=}}
\unsure{\operatorname{Brdt}}(C)/\operatorname{Pic}(A)
\simeq\sum_{\mathfrak p}\mathbb Z/e_{\mathfrak p}(C)\mathbb Z.
\]
\margin{$e_{\mathfrak p}(C)$: degré de ramification}

\deleted{Mais} Soit maintenant $A'$ une $A$-algèbre (comm.)
\deleted{Dans le cas}
et $C'$ une catégorie de Deuring sur $A'$, \inserted{et $T'\in V_{\unsure{A'|A}}(C')$,}
\unsure{supposons} $(C_{A'},T_{A'})\underset{A'}{\simeq}(C',\unsure{T'|A'})$.
\deleted{Alors l'ensemble des}
(il suffit pour ceci que $C_{A'}\simeq C'$ et $\operatorname{Pic}(A')=0$, p. ex. $A'$ local)
Alors l'ens. des équivalences de $(C_{A'},T_{A'})$ vers $(C',T')$, i.e. des couples $(\varphi,u)$ avec $\varphi$
\deleted{\unsure{équivalence}}
\diagram{\begin{tikzcd}
C_{\unsure{A'}} \arrow[r,"\varphi"] & C'
\end{tikzcd}}
\deleted{\unsure{de}} $A$-linéaire \unsure{admis.}, induisant $\simeq C_{A'}\longrightarrow C'$, et
\[
u:\nu(\varphi)(T_{A'})\xrightarrow{\sim}T',
\]
est une \uline{catégorie} \uline{torseur} sous la \inserted{Gr.-}catégorie
\[
\unsure{\mathrm{Eq}}\ \uline{\mathrm{Eq}}(C_{A'},T_{A'})
\quad\text{(à droite)}\quad
\text{\unsure{et}}\quad
\uline{\mathrm{Eq}}(C',T')\quad\text{(à gauche)}.
\]

Si p. ex. $\operatorname{Ram}(C'/A')=0$ ($C'$ non ramifiée sur $A'$)
\deleted{\unsure{on doit}} alors c'est une catégorie-torseur sous
\deleted{\unsure{$\mathcal E$}} $\uuline{\operatorname{Tors}}(S',\mu_n)$.
En fait, \unsure{on a} une équivalence \uline{canonique}
\[
\uline{\mathrm{Eq}}(C_{A'},T_{A'})
\simeq
\uline{\mathrm{Eq}}(C',T')
\]
(comme une \deleted{\unsure{Gr.-}} catégorie de \uuline{Picard}) et par suite les deux structures de tors.-catégories \unsure{peuvent} s'identifier \unsure{\ldots}
\margin{\uuline{À vérifier!!}}
